\documentclass[11pt]{article}
\usepackage[margin=1in]{geometry}
\usepackage{booktabs}
\usepackage{hyperref}

\title{From Reading Rules to Enforcing Actions:\\Where Solver-Checked Coding-Agent Policies Win}
\author{CS2520 Project}
\date{December 2026}

\begin{document}
\maketitle

% Topic sentences only (draft skeleton v2, 2026-11-24). New direction: runtime enforcement.
% Previous clash-only skeleton: main_v1_clash.tex. Numbers marked TBD are pending.

\begin{abstract}
Coding agents follow growing folders of memory rules, and we show that checking whether these
rules clash is a reading problem an LLM handles better than translation plus Z3, while
checking whether a concrete agent action breaks rules that were translated once and reviewed
is where a solver wins on [accuracy / cost / repeatability: TBD].
\end{abstract}

\section{Introduction}
Memory folders and \texttt{CLAUDE.md} files grow by accretion, so an agent must both cope with
rules that disagree and obey the rules on every git and file action it takes.

\section{Related Work}
\subsection{Agent Instruction Files}
Recent studies catalogue the content, smells, and divergences of real agent instruction files
and measure how agents behave when rules collide, but none compares a solver pipeline with a
whole-folder LLM on real memory folders.
\subsection{Solver-Aided Agent Policies}
Compile-then-enforce systems translate a natural-language policy into formal constraints
once and gate each tool call with a solver, and we adopt this structure from solver-aided
policy compliance work on tool-using agents.
\subsection{LLMs with Logic Solvers}
Hybrids of LLMs and solvers win when translation is easy and reasoning is deep, and tie or
lose on naturally written text, which matches what we find for clash detection.

\section{Two Jobs: Reading Rules and Enforcing Actions}
Finding clashes asks whether English rules can ever be followed together, while enforcement
asks whether one concrete, structured action breaks a fixed set of formal rules, and the
second removes most of the language understanding from each check.

\section{Method}
\subsection{Vocabulary}
A fixed vocabulary of git, file, and task actions with branch, path, and history conditions
lets in-scope rules be written as obligation, prohibition, permission, and conditional clauses.
\subsection{Compile Once: Translation and Human Review}
An LLM translates each rule in two steps, and a person reviews and fixes the translations once,
before any test action exists, with every edit counted.
\subsection{Enforce Every Action with Z3}
Each proposed tool call becomes a structured event by plain code, and Z3 decides whether it
breaks any compiled rule and names the rule and the situation that it breaks.
\subsection{Maintain: Checking New Rules for Clashes}
When a rule is added, the pair and group clash checker tests it against the compiled rules and
sends any clash to the person who reviews the update.

\section{Data}
We collected 1,287 rules from five real memory folders, translated 84 of them three times, and
built labelled clash sets and a set of agent actions per enforceable rule.

\section{Experiments and Results}
\subsection{Clash Detection: The LLM Reads Better}
On planted pairs, three-rule groups, and natural folders, one whole-folder LLM call finds as
many or more clashes than translation plus Z3, because translation and vocabulary coverage
(5 of 22 real clashes) lose most cases before the solver sees them.
\subsection{Runtime Enforcement: Setup}
We compare Z3 with automatic and with reviewed translations against Codex and DeepSeek judging
the same actions from the whole folder or from the selected rules only, under win conditions
fixed before any run.
\subsection{Runtime Enforcement: Results}
TBD: violation precision, recall, and F1, rule attribution, and agreement across repeats for
each arm.
\subsection{Cost and Latency}
TBD: one-time compile cost against the marginal cost and time per 100 and 1,000 checked actions.
\subsection{What the Human Review Adds}
TBD: the gap between automatic and reviewed translations, and how many translations the
reviewer changed.

\section{Discussion}
A solver is exact about what it receives, so it pays off when the hard reading is done once
and reused across many cheap checks, not when every question needs fresh reading of English.

\section{Limitations}
Enforcement covers only rules inside the vocabulary and cannot judge future obligations from one
action, the action set is small and model-written, the reviewed arm depends on human quality,
and our clash labels come from model labellers on five folders.

\section{Conclusion}
For coding-agent memory, use an LLM to read and find clashes, and use reviewed formal rules
with a solver to enforce git and file actions.

\end{document}
